% =============================================================================
% Chapter 11: Hardware Acceleration — ML Systems Lecture Slides
% =============================================================================
% Largest chapter in Vol1 (32 sections). Selective coverage focusing on:
%   - Why specialization matters (CPU vs GPU die area trade-off)
%   - Compute primitives (systolic arrays, tensor cores)
%   - The Memory Wall and Energy Hierarchy (Horowitz numbers)
%   - Roofline Model (arithmetic intensity, ridge points)
%   - Kernel fusion and dataflow optimization
%   - The accelerator spectrum (flexibility vs efficiency)
%
% IMAGES:
%   - cpu-vs-gpu-die.pdf        - memory-hierarchy-energy.pdf
%   - systolic-array-dataflow.pdf  - roofline-model.pdf
%   - kernel-fusion-before-after.pdf  - compute-bw-divergence.pdf
%   - accelerator-spectrum.pdf
% =============================================================================
\documentclass[aspectratio=169, 12pt]{beamer}
\usepackage{../../assets/beamerthememlsys}

\mlsyssetup{
  volume       = {Volume I},
  chapter      = {Chapter 11},
  logo         = {../../assets/img/logo-mlsysbook.png},
  instlogo     = {../../assets/img/logo-harvard.png},
  chaptertitle = {Hardware Acceleration},
}

% --- Fonts ---
\usepackage{fontspec}
\setsansfont{Helvetica Neue}[
  BoldFont={Helvetica Neue Bold},
  ItalicFont={Helvetica Neue Italic},
  BoldItalicFont={Helvetica Neue Bold Italic},
]
\setmonofont{JetBrains Mono}[Scale=0.85]

% --- Packages ---
\usepackage{booktabs}
\usepackage{amsmath}

% --- Image paths ---
\graphicspath{
  {images/}
}

% --- Chapter-specific macros ---
\newcommand{\AI}{\text{AI}}
\newcommand{\flopperbyte}{\text{FLOP/byte}}

% --- Helper: safe image include ---
\newcommand{\safeimg}[2][width=\textwidth,keepaspectratio]{%
  \IfFileExists{images/#2}{\includegraphics[#1]{#2}}{%
    \IfFileExists{#2}{\includegraphics[#1]{#2}}{%
      \fbox{\parbox[c][2.5cm][c]{0.85\linewidth}{\centering\footnotesize\textcolor{midgray}{[Missing image]}}}%
    }%
  }%
}

% --- Section count for navigation (must match actual \section{} count) ---
\setcounter{mlsystotalsections}{8}

\title{Hardware Acceleration}
\author{Vijay Janapa Reddi}
\institute{Harvard University}
\date{}

\begin{document}

% =============================================================================
% TITLE SLIDE
% =============================================================================
\mlsystitle{Hardware Acceleration}{Moving Data Costs More Than Computing It}{cover_ai_hardware.png}

% =============================================================================
% LEARNING OBJECTIVES
% =============================================================================
\begin{frame}{Learning Objectives}
\note{
% -- LINK: Ch10 compressed models; this chapter explains the silicon they run on
Model compression reduced what needs to be computed. This chapter reveals
the physical constraints of the hardware that executes those computations.

% -- NARRATE: Walk through objectives, emphasize the ``why'' over the ``what''
Read each objective. Stress: ``This chapter is not about which GPU is fastest.
It is about why hardware behaves the way it does --- energy, bandwidth,
arithmetic intensity. These are physics, not marketing.''

% -- ENGAGE: Surface the peak-TFLOPS fallacy early
Ask: ``How many of you have compared GPUs by peak TFLOPS?'' Most hands go up.
Say: ``By the end of this lecture, you will know why that is the wrong metric.''

% -- WARN: Students equate hardware knowledge with spec-sheet memorization
Common error: thinking this chapter is about memorizing GPU specs.
Correct framing: it teaches a diagnostic framework (Roofline) that applies
to any hardware, any workload, any generation.

% -- FLEX: [CORE] Learning objectives set the chapter roadmap
[CORE] This slide frames every subsequent concept.
IF AHEAD: ``Which objective surprises you most?''
IF SHORT: Read objectives quickly, skip the audience question.
}

\footnotesize
\begin{enumerate}\setlength\itemsep{0pt}
  \item Explain why \textbf{systolic arrays} and \textbf{tensor cores} achieve 10--100$\times$ better efficiency than general-purpose processors
  \item Calculate \textbf{arithmetic intensity} and use the \textbf{Roofline Model} to classify compute-bound vs.\ memory-bound workloads
  \item Predict bottlenecks by quantifying the \textbf{Memory Wall}: bandwidth, energy, and cache trade-offs
  \item Select \textbf{dataflow strategies} (weight-stationary, output-stationary) based on workload reuse
  \item Analyze \textbf{kernel fusion}, \textbf{tiling}, and \textbf{memory planning} for efficient execution
  \item Evaluate accelerator choices using \textbf{cost-performance} reasoning, not peak FLOPS
\end{enumerate}

\end{frame}

\begin{frame}{Visual Language}
\note{
% -- NARRATE: Brief orientation to the color code for this hardware chapter
Point to each card. In this chapter, blue highlights compute units (ALUs,
Tensor Cores), green highlights data paths and memory, orange highlights
scheduling/dataflow, and red highlights bottlenecks (memory wall, stalls).

% -- FLEX: [OPTIONAL] Skip if students saw this in Ch10
[OPTIONAL] If students already know the color system, say ``same legend'' and advance.
}

\small
Throughout this course, colors carry meaning:

\vspace{0.3cm}
\begin{columns}[T]
  \begin{column}{0.45\textwidth}
    \begin{mlsyscard}{computestroke}
    \textbf{Blue} --- Compute / Processing\\
    {\footnotesize GPU ops, forward/backward pass, inference}
    \end{mlsyscard}
    \vspace{0.15cm}
    \begin{mlsyscard}{datastroke}
    \textbf{Green} --- Data / Memory\\
    {\footnotesize Data flow, caches, healthy paths}
    \end{mlsyscard}
  \end{column}
  \begin{column}{0.45\textwidth}
    \begin{mlsyscard}{routingstroke}
    \textbf{Orange} --- Routing / Scheduling\\
    {\footnotesize Load balancers, batch windows}
    \end{mlsyscard}
    \vspace{0.15cm}
    \begin{mlsyscard}{errorstroke}
    \textbf{Red} --- Error / Cost / Bottleneck\\
    {\footnotesize Loss, decode phase, waste}
    \end{mlsyscard}
  \end{column}
\end{columns}
\end{frame}


% =============================================================================
\section{Why Specialize?}
% =============================================================================

\begin{frame}{The Central Surprise of Modern Computing}
\note{
% -- LINK: Ch10 compressed models to reduce data; this slide shows why data movement is the enemy
Model compression reduced bytes. This slide reveals the physical reason:
moving data costs 20,000x more energy than computing on it.

% -- NARRATE: State the thesis, then walk through the table and Iron Law column
Say: ``Here is the central surprise of modern computing. Arithmetic is nearly
free. Moving data is expensive. A modern accelerator computes 1,000 operations
in the time it takes to fetch one value from DRAM.'' Point to the Iron Law
table: ``Hardware acceleration increases the denominators --- peak throughput
and bandwidth. That is what this chapter is about.''

% -- ENGAGE: Pose the paradox
Ask: ``If computation is cheap, why do we need faster hardware?''
Give 20 seconds. Expected: because the bottleneck is data movement, not
arithmetic --- and faster hardware means better memory systems, not just
more ALUs.

% -- WARN: Students think ``faster hardware'' means more FLOPS
Common error: equating hardware acceleration with more arithmetic units.
Correct framing: the memory system (bandwidth, hierarchy, locality)
determines real performance far more than peak FLOPS.

% -- FLEX: [CORE] This is the chapter thesis --- everything else follows from it
[CORE] Every subsequent slide explains a consequence of this asymmetry.
IF AHEAD: ``What physical law makes data movement expensive?'' (capacitance, wire length)
IF SHORT: State the 20,000x number and the Iron Law connection, move on.
}

\footnotesize
\begin{columns}[T]
  \begin{column}{0.55\textwidth}
    \textbf{The Memory Wall}:
    \begin{itemize}\setlength\itemsep{0pt}
      \item In the time to fetch \textbf{one value} from DRAM, an accelerator performs \textbf{1,000+ operations}
      \item DRAM access costs \textbf{20,000$\times$} more energy than an INT8 add
      \item This is not engineering --- it is \textbf{physics}
    \end{itemize}

    \vspace{0.1cm}
    \begin{mlsyscard}{crimson}
    {\footnotesize The question is never ``which chip is fastest?'' but \alert{``which memory system best matches my access patterns?''}}
    \end{mlsyscard}
  \end{column}
  \begin{column}{0.42\textwidth}
    {\scriptsize
    \renewcommand{\arraystretch}{1.1}
    \begin{tabular}{@{}lr@{}}
      \toprule
      \textbf{Iron Law Term} & \textbf{Target} \\
      \midrule
      $D_{\text{vol}}/BW$    & Data selection \\
      $O/(R_p \cdot \eta)$   & Compression \\
      $R_p, BW, \eta$        & \textcolor{computestroke}{\textbf{HW Accel.}} \\
      $L_{\text{lat}}$       & Overhead \\
      \bottomrule
    \end{tabular}
    }

    \vspace{0.1cm}
    {\scriptsize Hardware acceleration increases the \textbf{denominators} $R_p$ and $BW$ in the Iron Law.}
  \end{column}
\end{columns}

\end{frame}

\begin{frame}{CPU vs GPU: The Silicon Area Trade-off}
\note{
% -- LINK: The Memory Wall showed data movement dominates; now show how chips allocate silicon
The previous slide stated the problem (data movement). This slide shows
two opposite architectural responses: CPUs optimize for single-thread
latency, GPUs optimize for aggregate throughput.

% -- NARRATE: Walk through die photos left to right
Point to CPU die: ``4 big cores, huge L3 cache, branch predictors, out-of-order
execution. Only about 5\% of the die is arithmetic.'' Point to GPU die:
``48+ streaming multiprocessors, tiny caches per SM. About 80\% of the die
is arithmetic.'' Point to FPGA: ``Reconfigurable fabric --- program the
hardware itself.'' Then: ``The GPU bet: sacrifice single-thread performance
for massive parallelism.''

% -- ENGAGE: What did GPUs sacrifice?
Ask: ``What did GPUs give up to get 80\% compute area?'' Give 20 seconds.
Expected: branch prediction, out-of-order execution, large caches,
single-thread performance. Each thread is slow, but there are thousands.

% -- WARN: Students think GPUs are universally faster than CPUs
Common error: ``GPUs are always faster because they have more FLOPS.''
Correct framing: GPUs are faster only for parallel, regular workloads.
For serial, branchy code (e.g., graph traversal), CPUs win.

% -- FLEX: [CORE] The CPU-GPU architectural trade-off is foundational
[CORE] This slide grounds every subsequent hardware discussion.
IF AHEAD: ``Where do TPUs sit on this spectrum?''
IF SHORT: Show die photos, state the 5\% vs 80\% numbers, move on.
}

% --- Layout: diagram top, real die photos below ---
\centering
\safeimg[width=0.88\textwidth,height=3.2cm,keepaspectratio]{cpu-vs-gpu-die.pdf}

\vspace{0.05cm}
\footnotesize
\begin{columns}[T]
  \begin{column}{0.30\textwidth}
    \centering
    \safeimg[width=\textwidth,height=1.6cm,keepaspectratio]{intel-cpu-die.jpg}

    \textcolor{computestroke}{\textbf{CPU}}: $\sim$5\% compute\\
    {\scriptsize Intel Core die. CC BY 2.0.}
  \end{column}
  \begin{column}{0.30\textwidth}
    \centering
    \safeimg[width=\textwidth,height=1.6cm,keepaspectratio]{gpu-die-stacked.jpg}

    \textcolor{datastroke}{\textbf{GPU}}: $\sim$80\% compute\\
    {\scriptsize NVIDIA Turing. CC BY 2.0.}
  \end{column}
  \begin{column}{0.30\textwidth}
    \centering
    \safeimg[width=\textwidth,height=1.6cm,keepaspectratio]{xilinx-fpga.jpg}

    \textcolor{routingstroke}{\textbf{FPGA}}: Reconfigurable\\
    {\scriptsize Xilinx Virtex. CC BY-SA 3.0.}
  \end{column}
\end{columns}

\end{frame}

\begin{frame}{Amdahl's Law: The Acceleration Ceiling}
\note{
% -- LINK: Die layouts showed GPUs have massive parallelism; Amdahl shows its limits
GPUs devote 80\% of silicon to compute. But Amdahl's Law says the serial
fraction of the workload caps the maximum benefit of all that parallelism.

% -- NARRATE: Walk through the equation, then the table
Point to the equation: ``p is the parallel fraction, S is the accelerator
speedup. The serial fraction (1-p) is the ceiling.'' Point to the table:
``ResNet-50 is 95\% parallel --- max speedup 20x. GPT-2 generation is 80\%
parallel --- max speedup only 5x. DLRM embedding is 60\% parallel ---
max speedup 2.5x.'' Then the red card: ``If data loading is 10\% of total
time, even infinite GPU speedup gives only 10x.''

% -- ENGAGE: Quick Amdahl's calculation
Ask: ``If p=0.95 and S=1000, what is the total speedup?'' Give 15 seconds.
Answer: 1/(0.05 + 0.95/1000) = 1/0.05095 = 19.6x. Nearly 20x, not 1000x.

% -- WARN: Students assume buying a faster GPU always helps proportionally
Common error: ``the H100 is 3x faster than A100, so my job runs 3x faster.''
Correct framing: only the parallel fraction benefits. Serial bottlenecks
(data loading, Python overhead, preprocessing) remain unchanged.

% -- FLEX: [CORE] Amdahl's Law is referenced throughout the rest of the course
[CORE] This equation appears again in Ch12 (Benchmarking) and Vol2.
IF AHEAD: ``How does Gustafson's Law offer a different perspective?''
IF SHORT: State the equation, give the ResNet example, move on.
}

\footnotesize
\begin{columns}[T]
  \begin{column}{0.50\textwidth}
    $$\text{Speedup} = \frac{1}{(1-p) + \frac{p}{S}}$$

    \vspace{0.05cm}
    {\scriptsize
    \begin{tabular}{@{}ll@{}}
      $p$ & Parallel fraction (matrix math) \\
      $S$ & Accelerator speedup (100--1000$\times$) \\
      $1\!-\!p$ & Serial fraction (data loading, Python) \\
    \end{tabular}
    }

    \vspace{0.1cm}
    \begin{mlsyscard}{errorstroke}
    {\scriptsize \textbf{Pitfall:} If data loading is 10\% ($p$=0.9), even $S\!=\!\infty$ yields only \textbf{10$\times$} total speedup.}
    \end{mlsyscard}
  \end{column}
  \begin{column}{0.46\textwidth}
    \renewcommand{\arraystretch}{1.1}
    {\scriptsize
    \begin{tabular}{@{}lrr@{}}
      \toprule
      \textbf{Workload} & \textbf{$p$} & \textbf{Max Speedup} \\
      \midrule
      ResNet-50       & 0.95 & 20$\times$ \\
      GPT-2 (gen.)    & 0.80 & 5$\times$ \\
      DLRM (embed)    & 0.60 & 2.5$\times$ \\
      MobileNet       & 0.70 & 3.3$\times$ \\
      \bottomrule
    \end{tabular}
    }

    \vspace{0.1cm}
    {\scriptsize\mlsysconcept{Key:}{ResNet scales better --- more time in parallelizable matrix math.}}
  \end{column}
\end{columns}

\end{frame}

% --- ACTIVE LEARNING 1: Predict ---
\begin{frame}{Predict: Where Is the Bottleneck?}
\note{
% -- LINK: Amdahl's showed serial bottlenecks; this reveals operation-level bottlenecks
Amdahl's Law capped total speedup. This shows that even within the parallel
portion, different operations hit different ceilings on the same GPU.

% -- NARRATE: Present the paradox, then wait
``MatMul: 280 TFLOPS. Softmax: 6 TFLOPS. Same GPU. Why?'' Wait 60 seconds.
Ask 2--3 students. Do NOT reveal the answer --- the Roofline Model explains it next.

% -- ENGAGE: Prediction before the Roofline reveal
Students who guess ``memory bandwidth'' are correct. Students who guess
``Softmax is harder math'' are wrong --- that is the learning moment.

% -- FLEX: [CORE] Active learning moment 1 of 3 --- primes the Roofline
[CORE] This prediction creates need-to-know for the Roofline Model.
IF SHORT: Reduce to 30 seconds, collect 1--2 hypotheses.
}

\centering
\vspace{0.8cm}
{\Large\bfseries Think--Write--Share}

\vspace{0.4cm}
{\large A transformer model runs on an H100 GPU.\\
\texttt{MatMul} layers hit 280 TFLOPS. \texttt{Softmax} layers hit 6 TFLOPS.\\[0.2cm]
\alert{Why does the same GPU give such different throughput\\for operations in the same model?}}

\vspace{0.4cm}
{\normalsize Write your hypothesis. \textcolor{midgray}{(60 seconds)}}

\end{frame}

% =============================================================================
\section{Compute Primitives}
% =============================================================================

\begin{frame}{The MAC: Neural Networks' Dominant Operation}
\note{
% -- LINK: Die layouts showed 80\% compute on GPUs; now show what that compute does
GPUs devote 80\% of silicon to arithmetic. This slide reveals what dominates:
multiply-accumulate (MAC) operations, 95\% of execution time.

% -- NARRATE: Walk through the primitives table and code example
``Three primitives: vector (ReLU), matrix (Dense, Conv2D), and special functions
(Softmax).'' Point to code: ``A dense layer is just three nested loops of
multiply-accumulate. 32 x 512 x 256 = 4.2M MACs.'' Key: their regularity
makes hardware specialization possible.

% -- ENGAGE: Why does regularity matter for hardware?
Ask: ``Why does the regularity of MACs enable hardware specialization?''
Expected: predictable patterns allow fixed-function pipelines and data reuse.

% -- WARN: Students may think all operations benefit equally from GPUs
Common error: assuming Softmax benefits as much as MatMul from GPU parallelism.
Correct framing: Softmax is memory-bound (low AI), MatMul is compute-bound.

% -- FLEX: [CORE] Establishes why systolic arrays and Tensor Cores exist
[CORE] This slide explains what the next two slides accelerate.
IF AHEAD: ``What about operations with branching (e.g., conditional layers)?''
IF SHORT: State ``MACs = 95\% of time'' and move to systolic arrays.
}

\footnotesize
\begin{columns}[T]
  \begin{column}{0.50\textwidth}
    \textbf{Three compute primitives} dominate all accelerators:

    \vspace{0.1cm}
    \renewcommand{\arraystretch}{1.1}
    {\scriptsize
    \begin{tabular}{@{}lll@{}}
      \toprule
      \textbf{Primitive} & \textbf{Example} & \textbf{Hardware} \\
      \midrule
      Vector ops  & ReLU, bias add   & SIMD units \\
      Matrix ops  & Dense, Conv2D    & Tensor Cores \\
      Special fn  & Softmax, GELU    & SFUs \\
      \bottomrule
    \end{tabular}
    }

    \vspace{0.1cm}
    \mlsysconcept{Key:}{MACs are 95\%+ of execution time.}
  \end{column}
  \begin{column}{0.46\textwidth}
    \begin{mlsyscard}{computestroke}
    {\scriptsize \textbf{Dense layer (256$\to$512):}\\[0.05cm]
    \texttt{for n in batch:}\\
    \texttt{~~for m in output:}\\
    \texttt{~~~~sum = bias[m]}\\
    \texttt{~~~~for k in input:}\\
    \texttt{~~~~~~sum += x[n,k]*w[k,m]}\\[0.05cm]
    $32 \times 512 \times 256$ = \textbf{4.2M MACs}
    }
    \end{mlsyscard}
  \end{column}
\end{columns}

\end{frame}

\begin{frame}{Systolic Arrays: Data Reuse at Scale}
\note{
% -- LINK: MACs dominate; systolic arrays are hardware designed specifically for MACs
The previous slide showed MACs are 95\% of execution. Systolic arrays are
silicon designed to execute MACs with maximum data reuse.

% -- NARRATE: Walk through the systolic array diagram
``Inputs flow right, weights flow down, partial sums accumulate in each PE.
The key: each weight is loaded once and used 256 times in TPUv1's 256x256
array. That is a 256x reduction in memory bandwidth demand.''
ANALOGY: ``Named after the heart's rhythmic pumping --- data pulses through
the array in a regular cadence.''

% -- ENGAGE: Calculate the bandwidth savings
Ask: ``If each weight is used 256 times, how much less bandwidth does the
systolic array need vs a naive implementation?'' Expected: 256x less.

% -- WARN: Students confuse systolic arrays with GPU cores
Common error: thinking systolic arrays are just ``more parallel cores.''
Correct framing: systolic arrays are a fundamentally different architecture ---
nearest-neighbor data passing, no shared memory, fixed dataflow.

% -- FLEX: [CORE] Systolic arrays are the foundation of TPU architecture
[CORE] Students need this to understand TPU vs GPU trade-offs.
IF AHEAD: ``What happens when the matrix is smaller than 256x256?''
IF SHORT: State the 256x reuse number, skip the analogy.
}

% --- Full-width diagram ---
\centering
\safeimg[width=0.88\textwidth,height=4.5cm,keepaspectratio]{systolic-array-dataflow.pdf}

\vspace{0.1cm}
\mlsysinsight{Data reuse:}{TPUv1's 256$\times$256 array = 65,536 MACs. Each weight loaded once, used 256 times.}

\end{frame}

\begin{frame}{Tensor Cores: GPU's Matrix Engine}
\note{
% -- LINK: Systolic arrays are TPU-specific; Tensor Cores are NVIDIA's equivalent
Systolic arrays were Google's approach. NVIDIA added Tensor Cores to GPUs
starting with Volta (2017) to get similar specialization within a GPU.

% -- NARRATE: Walk through the specs table
``Each Tensor Core: 4x4x4 matrix MAC per clock. A100: 432 Tensor Cores = 312
TFLOPS FP16. H100: 528 Tensor Cores = 989 TFLOPS FP16.'' Then the green card:
``1,000x growth in a decade --- from K20X to H100 --- through specialization,
not clock speed.'' Requirements: ``FP16/BF16, dimensions aligned to 8/16.''

% -- WARN: Students assume Tensor Cores are always active
Common error: ``My GPU has 989 TFLOPS.'' Correct framing: Tensor Cores
activate only with aligned FP16/BF16 data and specific matrix dimensions.
Unaligned or FP32 code runs on CUDA cores at much lower throughput.

% -- FLEX: [OPTIONAL] Can merge with systolic arrays if short
[OPTIONAL] The V100/A100/H100 progression is the key takeaway.
IF AHEAD: ``How do Sparse Tensor Cores handle 2:4 sparsity from Ch10?''
IF SHORT: Skip this slide, merge key numbers into systolic array discussion.
}

\footnotesize
\begin{columns}[T]
  \begin{column}{0.55\textwidth}
    \textbf{NVIDIA Tensor Core} (since Volta, 2017):
    \begin{itemize}\setlength\itemsep{0pt}
      \item Each core: $4\!\times\!4\!\times\!4$ matrix MAC per clock
      \item A100: \textbf{432 Tensor Cores} $\to$ 312 TF (FP16)
      \item H100: \textbf{528 Tensor Cores} $\to$ 989 TF (FP16)
    \end{itemize}

    \vspace{0.05cm}
    \textbf{Peak utilization requires:}
    {\scriptsize FP16/BF16 inputs, dimensions aligned to 8/16, sufficient batch size.}
  \end{column}
  \begin{column}{0.42\textwidth}
    \renewcommand{\arraystretch}{1.1}
    {\scriptsize
    \begin{tabular}{@{}lrr@{}}
      \toprule
      \textbf{GPU} & \textbf{Tensor Cores} & \textbf{FP16 TF} \\
      \midrule
      V100 & 640  & 125 \\
      A100 & 432  & 312 \\
      H100 & 528  & 989 \\
      \bottomrule
    \end{tabular}
    }

    \vspace{0.1cm}
    \begin{mlsyscard}{datastroke}
    {\scriptsize \textbf{1,000$\times$} growth in one decade (K20X$\to$H100) --- specialization, not clock speed.}
    \end{mlsyscard}
  \end{column}
\end{columns}

\end{frame}

% =============================================================================
\section{The Memory Wall}
% =============================================================================

\mlsysfocus{The Memory Wall}{%
\small Moving data costs \textbf{20,000$\times$} more energy than computing it.\\[0.3cm]
\normalsize $\underbrace{\text{INT8 Add}}_{\text{0.03 pJ}} \quad\longleftrightarrow\quad \underbrace{\text{DRAM Read}}_{\text{640 pJ}}$\\[0.3cm]
\small This is physics, not engineering.%
}

\begin{frame}{The Horowitz Numbers: Energy Per Operation}
\note{
% -- LINK: The focus frame stated 20,000x dramatically; now provide the precise numbers
The dark-background slide made the emotional impact. This slide provides
the exact Horowitz (ISSCC 2014) energy numbers that every hardware
designer memorizes.

% -- NARRATE: Walk through the table row by row, emphasizing the 200x gap
``An FP32 multiply costs 3.7 pJ. A DRAM read costs 640 pJ --- that is
200$\times$ more energy to \emph{fetch} a number than to \emph{multiply}
it.'' Point to the ratio column: ``INT8 add is the cheapest at 0.03 pJ.
DRAM is 20,000$\times$ more. This is why quantization from Ch.\ 10 saves
energy: INT8 values are 4$\times$ smaller, so 4$\times$ fewer DRAM reads.''

% -- ENGAGE: Calculate total energy for one dense layer
Ask: ``A dense layer reads 1M weights from DRAM and does 1M MACs. What
fraction of energy is compute vs.\ data movement?'' Give 20 seconds.
Expected: Compute = 1M $\times$ 3.7 pJ = 3.7 $\mu$J. DRAM = 1M $\times$ 640 pJ
= 640 $\mu$J. Data movement is 99.4\% of total energy.

% -- WARN: Students focus on compute speed, not energy cost
Common error: ``FLOPs determine energy.'' Correct framing: data movement
dominates energy by 100--200$\times$. Optimizing compute is irrelevant
if you do not reduce data movement.

% -- FLEX: [CORE] These numbers are the physical foundation for everything
[CORE] Referenced by fusion, tiling, quantization, and serving chapters.
IF AHEAD: ``How have these numbers changed since 2014 with HBM3?''
IF SHORT: State the 200$\times$ ratio and the 99.4\% fraction.
}

\footnotesize
\begin{columns}[T]
  \begin{column}{0.55\textwidth}
    \textbf{Horowitz, ISSCC 2014} (45 nm reference):

    \vspace{0.1cm}
    \renewcommand{\arraystretch}{1.05}
    {\scriptsize
    \begin{tabular}{@{}llrr@{}}
      \toprule
      \textbf{Operation} & \textbf{Width} & \textbf{Energy (pJ)} & \textbf{Ratio} \\
      \midrule
      INT8 Add      & 8-bit   & 0.03    & 1$\times$ \\
      INT8 Multiply & 8-bit   & 0.2     & 7$\times$ \\
      FP16 Multiply & 16-bit  & 1.1     & 37$\times$ \\
      FP32 Multiply & 32-bit  & 3.7     & 123$\times$ \\
      SRAM Read     & 32-bit  & 5       & 167$\times$ \\
      \textcolor{errorstroke}{\textbf{DRAM Read}} & 64-bit & \textcolor{errorstroke}{\textbf{640}} & \textcolor{errorstroke}{\textbf{20,000$\times$}} \\
      \bottomrule
    \end{tabular}
    }
  \end{column}
  \begin{column}{0.42\textwidth}
    \vspace{0.1cm}
    \begin{mlsyscard}{errorstroke}
    {\scriptsize \textbf{Dense layer energy split:}\\[0.05cm]
    1M MACs $\times$ 3.7 pJ = 3.7 $\mu$J\\
    1M DRAM reads $\times$ 640 pJ = 640 $\mu$J\\[0.05cm]
    Data movement: \textbf{99.4\%}\\
    Compute: \textbf{0.6\%}}
    \end{mlsyscard}

    \vspace{0.1cm}
    \mlsysconcept{Implication:}{Reducing data movement by 2$\times$ saves more energy than making ALUs 100$\times$ faster.}
  \end{column}
\end{columns}
\mlsyscite{M. Horowitz, ISSCC 2014}

\end{frame}

\begin{frame}{The Energy Hierarchy (Horowitz Numbers)}
\note{
% -- LINK: The focus frame stated 20,000x; this slide shows the full hierarchy
The dark-background slide stated the 20,000x number dramatically. This slide
provides the complete energy pyramid from registers to host DRAM.

% -- NARRATE: Walk through the pyramid from bottom to top
``Registers: fastest, smallest, cheapest per access. L1 SRAM: 5x more energy.
L2: 10x. HBM/DRAM: 128x vs SRAM, 20,000x vs INT8 add.'' Then: ``Every level
trades capacity for speed and energy. This is why all accelerator design is
fundamentally about data locality.''

% -- ENGAGE: What does this mean for neural network layer design?
Ask: ``What does this hierarchy mean for how we design neural network layers?''
Expected: layers should maximize data reuse (high arithmetic intensity),
keep working sets in SRAM, minimize DRAM traffic.

% -- WARN: Students focus on compute speed, not energy cost
Common error: ``My GPU is fast enough.'' Correct framing: even if compute is
fast, DRAM access dominates energy. Battery life and thermal limits are set
by data movement, not arithmetic.

% -- FLEX: [CORE] The energy hierarchy explains every hardware design decision
[CORE] This is the physical foundation for fusion, tiling, and dataflow.
IF AHEAD: ``How does HBM3 change these numbers vs DDR5?''
IF SHORT: State the 20,000x and 128x numbers, move on.
}

% --- Full-width diagram ---
\centering
\safeimg[width=0.88\textwidth,height=4.5cm,keepaspectratio]{memory-hierarchy-energy.pdf}

\vspace{0.1cm}
\mlsysalert{Design implication:}{Accelerators must prioritize \emph{data locality} over raw arithmetic throughput.}

\end{frame}

\begin{frame}{The Compute-Bandwidth Divergence}
\note{
% -- LINK: Energy hierarchy was static; this shows the trend is getting worse
The Horowitz numbers showed the energy gap. This slide shows the gap is
widening: compute grows faster than bandwidth every generation.

% -- NARRATE: Walk through the divergence diagram
``Compute: 8x growth V100 to H100. Bandwidth: only 3.7x. The ridge point ---
operations per byte needed to saturate the chip --- has doubled.'' Then the
counterintuitive result: ``Code that was compute-bound on A100 may become
memory-bound on H100. Upgrading hardware can shift your bottleneck regime.''

% -- ENGAGE: What happens to your kernel when you upgrade GPUs?
Ask: ``If your kernel achieves 80\% utilization on A100, what happens on H100?''
Expected: may drop to 40\% if it was near the ridge point, because H100's
ridge point is nearly 2x higher.

% -- WARN: Students assume new hardware always helps
Common error: ``H100 is 3x faster, so my workload runs 3x faster.''
Correct framing: only if your workload is compute-bound. Memory-bound
workloads see gains only proportional to bandwidth improvement (1.7x).

% -- FLEX: [CORE] The divergence trend is the most important hardware trend
[CORE] This explains why kernel fusion and quantization matter more each year.
IF AHEAD: ``When will bandwidth catch up?'' (Never --- physics of wires vs transistors.)
IF SHORT: State the 8x vs 3.7x numbers and the doubled ridge point.
}

% --- Full-width diagram ---
\centering
\safeimg[width=0.88\textwidth,height=4.5cm,keepaspectratio]{compute-bw-divergence.pdf}

\vspace{0.1cm}
\small
\begin{columns}[T]
  \begin{column}{0.45\textwidth}
    \centering
    \textcolor{computestroke}{\textbf{Compute:}} 8$\times$ growth (V100$\to$H100)
  \end{column}
  \begin{column}{0.45\textwidth}
    \centering
    \textcolor{routingstroke}{\textbf{Bandwidth:}} 3.7$\times$ growth (V100$\to$H100)
  \end{column}
\end{columns}
\vspace{0.05cm}
{\small The ridge point doubled: keeping hardware fully utilized gets \textbf{harder} each generation.}

\end{frame}

% =============================================================================
\section{Roofline Model}
% =============================================================================

\begin{frame}{Arithmetic Intensity: The Key Metric}
\note{
% -- LINK: Divergence showed the gap widens; AI is the metric that diagnoses it
The divergence slide showed compute outpacing bandwidth. Arithmetic intensity
is the single number that determines which ceiling you hit.

% -- NARRATE: Walk through equation, then the operations table
``AI = FLOPs / Bytes Transferred. Low AI = memory-bound. High AI =
compute-bound. The ridge point is where they meet.'' A100 ridge: 156.
H100 ridge: 296. Point to table: ``Conv2D and Dense MatMul are compute-bound.
Everything else --- Softmax, LayerNorm, ReLU, Embedding --- is memory-bound.
Most operations are memory-bound!''

% -- ENGAGE: Revisit the prediction exercise
Say: ``Now you can answer the prediction: MatMul has AI of 64--256 (above
ridge). Softmax has AI of 2--5 (far below). Same GPU, different ceilings.''

% -- WARN: Students assume high-FLOP operations are compute-bound
Common error: ``Softmax does a lot of math, so it must be compute-bound.''
Correct framing: Softmax reads/writes the full tensor for each of exp, sum,
divide --- the bytes dominate, not the FLOPs.

% -- FLEX: [CORE] AI is the most important diagnostic metric in the course
[CORE] Students use this in the exercise next and in Ch12 (Benchmarking).
IF AHEAD: ``How does batching change the AI of a dense layer?''
IF SHORT: State the equation, show the table, highlight ``most ops memory-bound.''
}

\footnotesize
\begin{columns}[T]
  \begin{column}{0.48\textwidth}
    $$\text{AI} = \frac{\text{FLOPs}}{\text{Bytes Transferred}}$$

    \vspace{0.05cm}
    \textbf{Which ceiling do you hit?}
    \begin{itemize}\setlength\itemsep{0pt}
      \item Low AI $\to$ \textcolor{errorstroke}{\textbf{memory-bound}}
      \item High AI $\to$ \textcolor{computestroke}{\textbf{compute-bound}}
    \end{itemize}

    \vspace{0.05cm}
    \textbf{Ridge Point} = $R_{\text{peak}} / BW$

    \vspace{0.05cm}
    \begin{mlsyscard}{crimson}
    {\scriptsize \textbf{A100:} $312\text{ TF}/2\text{ TB/s}=\textbf{156}$ FLOP/byte\\
    \textbf{H100:} $989\text{ TF}/3.35\text{ TB/s}=\textbf{296}$ FLOP/byte}
    \end{mlsyscard}
  \end{column}
  \begin{column}{0.48\textwidth}
    \renewcommand{\arraystretch}{1.05}
    {\scriptsize
    \begin{tabular}{@{}lrl@{}}
      \toprule
      \textbf{Operation} & \textbf{AI} & \textbf{Regime} \\
      \midrule
      Conv2D        & 50--200 & Compute \\
      Dense MatMul  & 64--256 & Compute \\
      Depthwise Conv & 10--20 & Memory \\
      Softmax       & 2--5   & Memory \\
      LayerNorm     & 5--10  & Memory \\
      Embedding     & $<$1   & Memory \\
      ReLU          & 0.125  & Memory \\
      \bottomrule
    \end{tabular}
    }

    \vspace{0.05cm}
    {\footnotesize\textcolor{errorstroke}{\textbf{Most operations are memory-bound!}}}
  \end{column}
\end{columns}

\end{frame}

\begin{frame}{The Roofline Plot}
\note{
% -- LINK: AI defined the metric; the Roofline Plot visualizes it
Arithmetic intensity was the number. The Roofline Plot is the visual tool
that maps every operation to its performance ceiling.

% -- NARRATE: Walk through the diagram's two ceilings
``Two lines: the sloped bandwidth ceiling (left) and the flat compute ceiling
(right). They meet at the ridge point.'' Point to where each ML operation sits:
``Most cluster to the LEFT of the ridge --- memory-bound. ReLU is 2,400x below
the H100 roofline.'' Then: ``This is why kernel fusion is the most important
optimization --- it moves operations rightward by combining their AI.''

% -- ENGAGE: Where does your favorite model sit?
Ask: ``Based on the AI table, where does transformer inference at batch=1
sit on this plot?'' Expected: far left, deeply memory-bound.

% -- WARN: Students think the Roofline is only theoretical
Common error: ``This is a model, not reality.'' Correct framing: the Roofline
is validated by profiling tools (NVIDIA Nsight, rocprof). Real kernels land
within 10--20\% of the predicted ceiling.

% -- FLEX: [CORE] The Roofline is THE diagnostic tool of the course
[CORE] Referenced in Ch12 (Benchmarking) and Vol2.
IF AHEAD: ``How do caches create multiple rooflines?''
IF SHORT: Show the plot, point to the ridge, state ``most ops are left.''
}

% --- Full-width diagram ---
\centering
\safeimg[width=0.92\textwidth,height=4.8cm,keepaspectratio]{roofline-model.pdf}

\vspace{0.1cm}
\mlsysinsight{Diagnostic power:}{When inference is slow, the Roofline tells you whether to optimize compute, memory, or software.}

\end{frame}

\begin{frame}{Quick Check}
\note{
% -- NARRATE: Pause for retrieval before the exercise
``Close notes. What does the ridge point of the Roofline Model represent?''
Wait 30 seconds. Expected: the arithmetic intensity at which the workload
transitions from memory-bound to compute-bound (R_peak / BW).

% -- FLEX: [CORE] Distributed retrieval cements the Roofline before the exercise
[CORE] 30-second pause between concept and application.
IF SHORT: Reduce to 15 seconds, still require an answer.
}

\centering
\Large\bfseries Quick check --- no peeking.\\[0.5cm]
\normalsize What does the ridge point of the Roofline Model represent?\\[0.3cm]
{\small 30 seconds --- then we continue.}
\end{frame}



% --- ACTIVE LEARNING 2: Exercise ---
\begin{frame}{Your Turn: Roofline Diagnosis}
\note{
% -- LINK: The Roofline Plot was introduced; now apply it to a concrete layer
Students saw the Roofline visually. This exercise forces them to calculate
arithmetic intensity and diagnose a real layer.

% -- NARRATE: Present problem, then reveal after pause
``Dense layer: 2048 to 512, FP16, batch=1. Calculate AI.'' Wait 90 seconds.
Solution: ``AI = 2.1M FLOPs / 2.1M bytes = 1 FLOP/byte. Ridge is 156.
Deeply memory-bound! GPU achieves under 1\% utilization. Fix: increase batch
size. At batch=64: AI = 60, still below ridge but 40x better.''

% -- ENGAGE: Peer instruction protocol
Present (30s). Individual work (60s). If 30--70\% correct: pair discussion
(90s), re-vote. The key insight: batch size transforms the regime.

% -- WARN: Students forget to include all memory traffic
Common error: counting only weight bytes, forgetting input and output tensors.
Correct framing: AI denominator includes ALL bytes transferred --- inputs,
weights, and outputs.

% -- FLEX: [CORE] Active learning moment 2 of 3
[CORE] Central calculation that connects AI to hardware utilization.
IF AHEAD: ``At what batch size does this layer cross the A100 ridge?''
IF SHORT: Show solution immediately, walk through the math.
}

\footnotesize
\begin{columns}[T]
  \begin{column}{0.55\textwidth}
    {\normalsize\bfseries Diagnose This Layer}

    \vspace{0.15cm}
    Dense layer: 2048 inputs $\to$ 512 outputs, FP16.

    \vspace{0.1cm}
    \textbf{At batch = 1:}
    \begin{itemize}\setlength\itemsep{0pt}
      \item FLOPs = $2 \times 2048 \times 512$ = 2.1M
      \item Bytes = $(2048\!\times\!512 + 2048 + 512)\!\times\!2$ = 2.1 MB
    \end{itemize}

    \vspace{0.05cm}
    \textbf{Calculate} the arithmetic intensity.\\
    Compute-bound or memory-bound on A100 (ridge = 156)?

    {\scriptsize\textcolor{midgray}{(90 seconds --- compare with a neighbor)}}
  \end{column}
  \begin{column}{0.42\textwidth}
    \pause
    \begin{mlsyscard}{datastroke}
    {\scriptsize \textbf{Solution:}\\[0.05cm]
    AI = $2.1\text{M}/2.1\text{M}$ = \textbf{1 FLOP/byte}\\[0.05cm]
    Ridge = 156 $\to$ \textcolor{errorstroke}{\textbf{Deeply memory-bound!}}\\[0.05cm]
    GPU achieves $<$1\% utilization.\\[0.05cm]
    \textbf{Fix:} Increase batch size.\\
    At batch=64: AI $\approx$ 60 (40$\times$ better)
    }
    \end{mlsyscard}
  \end{column}
\end{columns}

\end{frame}

% =============================================================================
\section{Dataflow \& Fusion}
% =============================================================================

\begin{frame}{Dataflow Strategies: Which Data Stays Local?}
\note{
% -- LINK: The Roofline showed memory-bound ops dominate; dataflow minimizes traffic
The Roofline revealed most ops are memory-bound. Dataflow strategies determine
which tensor stays in fast local memory to minimize DRAM traffic.

% -- NARRATE: Walk through three strategies in the table
``Weight-stationary: weights stay in registers, inputs/outputs stream through.
Best for CNNs and FC layers where weights are reused across the batch. TPU
uses this.'' ``Output-stationary: partial sums stay local, minimizing write
traffic.'' ``Input-stationary: inputs cached, good for small-batch inference.''
Then the red card: ``Data movement costs 100--1,000x more than compute.''

% -- ENGAGE: Which strategy for LLM autoregressive generation?
Ask: ``For autoregressive LLM generation at batch=1, which strategy?''
Expected: weight-stationary --- weights are huge and reused across tokens.

% -- WARN: Students think one dataflow is always best
Common error: ``Weight-stationary is always optimal.'' Correct framing: the
best strategy depends on which tensor is largest relative to local storage.

% -- FLEX: [CORE] Dataflow connects to systolic array design
[CORE] Explains why TPU chose weight-stationary.
IF AHEAD: ``Can you dynamically switch dataflow per layer?''
IF SHORT: Cover weight-stationary only, skip the other two.
}

\scriptsize
\renewcommand{\arraystretch}{1.05}
\begin{tabular}{@{}lp{3.2cm}p{3cm}p{3.2cm}@{}}
  \toprule
  & \textbf{Weight-Stationary} & \textbf{Output-Stationary} & \textbf{Input-Stationary} \\
  \midrule
  \textbf{Local data}  & Weights in registers & Partial sums local & Inputs cached \\
  \textbf{From DRAM} & Inputs \& outputs & Weights \& inputs & Weights \& outputs \\
  \textbf{Best for}    & CNN, FC layers (weights reused across batch) & Large reductions (min write traffic) & Small-batch inference (inputs reused) \\
  \textbf{Example}  & TPU systolic array & NVDLA & Eyeriss \\
  \bottomrule
\end{tabular}

\vspace{0.1cm}
\begin{mlsyscard}{errorstroke}
{\scriptsize \textbf{Physical constraint:} Data movement costs 100--1,000$\times$ more energy than compute. Optimize for minimal data movement, not minimal FLOPs.}
\end{mlsyscard}

\end{frame}

\begin{frame}{Kernel Fusion: Eliminating Intermediate Writes}
\note{
% -- LINK: Dataflow keeps data local within one op; fusion keeps data local across ops
Dataflow strategies optimize within a single operation. Kernel fusion
eliminates intermediate DRAM writes between consecutive operations.

% -- NARRATE: Walk through the before/after diagram
Left: ``Conv2D writes to HBM. BatchNorm reads from HBM, writes again. ReLU
reads again. Six round-trips.'' Right: ``One fused kernel: one read, compute
all three in SRAM, one write. ResNet-50: 47 ms to 8 ms --- 6x speedup from
fusion alone.'' Then: ``Fusion converts memory-bound ops into one compute-bound
kernel by increasing their combined arithmetic intensity.''

% -- ENGAGE: Why is fusion the most impactful compiler optimization?
Ask: ``If most ops are memory-bound, why does fusing them help?'' Expected:
fusing N memory-bound ops eliminates N-1 intermediate writes, raising the
combined AI above the ridge point.

% -- WARN: Students think fusion only saves a little time
Common error: ``Fusion saves just the write time.'' Correct framing: the
intermediate writes are the dominant cost for memory-bound ops. Eliminating
them can give 3--6x speedup, not just a few percent.

% -- FLEX: [CORE] Fusion is the single most important compiler optimization
[CORE] This is the practical answer to the Memory Wall.
IF AHEAD: ``What limits fusion? When can't you fuse?''
IF SHORT: State the ResNet-50 number (47 to 8 ms), move on.
}

% --- Full-width diagram ---
\centering
\safeimg[width=0.88\textwidth,height=4.5cm,keepaspectratio]{kernel-fusion-before-after.pdf}

\vspace{0.1cm}
\mlsysconcept{Why it works:}{Fusion converts memory-bound ops (low AI individually) into one compute-bound kernel (high AI combined).}

\end{frame}

\begin{frame}{Tiling: Making Data Fit in Fast Memory}
\note{
% -- LINK: Fusion kept data local across ops; tiling keeps data local within one op
Fusion eliminated intermediate writes. Tiling ensures each individual
operation's data fits in fast SRAM during execution.

% -- NARRATE: Walk through the tiling process and AI impact
``Large matrices (4096x4096) don't fit in 128 KB SRAM. Tiling partitions
them into blocks that do fit. Process all computations for each tile before
moving on.'' Point to the blue card: ``Naive matmul: AI = 1 FLOP/byte
(reload every time). Tiled matmul: AI = 100+ FLOP/byte (reuse from SRAM).
Same math, 100x higher AI!''

% -- WARN: Students confuse tiling with batching
Common error: ``Tiling is just processing in smaller batches.'' Correct
framing: tiling partitions the matrix dimensions, not the batch dimension.
It maximizes reuse of operands already in SRAM.

% -- FLEX: [OPTIONAL] Tiling is important but secondary to fusion
[OPTIONAL] The AI comparison (1 vs 100+) is the key takeaway.
IF AHEAD: ``How does the compiler choose the optimal tile size?''
IF SHORT: State the 1 vs 100+ AI contrast, move on.
}

\small
\begin{columns}[T]
  \begin{column}{0.55\textwidth}
    \textbf{The problem:} Large matrices (e.g., $4096\times4096$) do not fit in fast SRAM (128 KB).

    \vspace{0.15cm}
    \textbf{The solution:} Partition into tiles that fit:
    \begin{itemize}\setlength\itemsep{0pt}
      \item Load tile of A, tile of B into SRAM
      \item Compute \emph{all} operations for this tile
      \item Write result tile back to DRAM
      \item Advance to next tile
    \end{itemize}

    \vspace{0.1cm}
    \mlsysconcept{Key:}{Process all computations per tile before moving on --- never bounce between tiles.}
  \end{column}
  \begin{column}{0.42\textwidth}
    \begin{mlsyscard}{computestroke}
    \textbf{Impact on Arithmetic Intensity:}\\[0.1cm]
    {\footnotesize
    Naive matmul: AI $\approx$ 1 FLOP/byte\\
    {\footnotesize (reload operands every time)}\\[0.1cm]
    Tiled matmul: AI $\approx$ 100+ FLOP/byte\\
    {\footnotesize (reuse data from SRAM)}\\[0.1cm]
    \textcolor{datastroke}{\textbf{Same math, 100$\times$ higher AI!}}
    }
    \end{mlsyscard}
  \end{column}
\end{columns}

\end{frame}

% --- ACTIVE LEARNING 3: Discussion ---
\begin{frame}{Discussion: Fusion or Faster Hardware?}
\note{
% -- LINK: Students learned fusion and Roofline; now synthesize them
Students know the Roofline (most ops memory-bound) and fusion (eliminates
intermediate writes). This forces them to choose between hardware and software.

% -- NARRATE: Frame the debate, then facilitate
``Your company offers two upgrades: A100 to H100 (3x TFLOPS), or hire a
compiler engineer to fuse memory-bound kernels. Which gives bigger speedup
for a transformer?'' 90 seconds pair discussion. Cold-call 2--3 pairs.
Answer: fusion is almost always more impactful because most ops are memory-bound.
A 2x faster GPU with the same bandwidth barely helps memory-bound ops.

% -- ENGAGE: Force a choice with reasoning
Students must commit to A or B and explain using the Roofline. The best
answers reference specific AI values and the ridge point.

% -- WARN: Students default to ``buy faster hardware''
Common error: hardware upgrades feel more concrete than compiler work.
Correct framing: 3x more TFLOPS helps only compute-bound ops. For memory-bound
Softmax, LayerNorm, etc., the upgrade is nearly invisible.

% -- FLEX: [CORE] Active learning moment 3 of 3 --- synthesis
[CORE] This integrates Roofline, fusion, and hardware selection.
IF AHEAD: ``What about buying a GPU with more bandwidth instead?''
IF SHORT: Reduce to 60 seconds, cold-call one pair.
}

\centering
\vspace{0.6cm}
{\large\bfseries Turn and Talk \textcolor{midgray}{(90 seconds)}}

\vspace{0.4cm}
{\large Your company offers two upgrades:\\[0.15cm]
\textbf{A:} Upgrade A100 $\to$ H100 (3$\times$ more TFLOPS)\\
\textbf{B:} Hire a compiler engineer to fuse memory-bound kernels\\[0.2cm]
\alert{Which gives a bigger speedup for a transformer?}\\[0.15cm]
{\normalsize \textit{Hint: Where do most ops sit on the Roofline?}}}

\end{frame}

% =============================================================================
\section{Accelerator Spectrum}
% =============================================================================

\begin{frame}{The Flexibility-Efficiency Trade-off}
\note{
% -- LINK: Fusion showed software optimization; now show the hardware design spectrum
Software (fusion, tiling) optimizes execution on fixed hardware. This slide
shows how different hardware designs trade flexibility for efficiency.

% -- NARRATE: Walk through the spectrum left to right
``CPU: most flexible, least efficient. GPU: flexible enough for training,
efficient enough for most workloads. TPU: sacrifices flexibility for higher
efficiency at matrix ops. FPGA: reconfigurable but slower. ASIC/NPU: fastest
for fixed workloads, zero flexibility.'' Then the green insight: ``Match the
accelerator's memory system to your workload's access patterns.''

% -- WARN: Students think more specialization is always better
Common error: ``ASICs are fastest, so always use ASICs.'' Correct framing:
ASICs cannot adapt to new architectures. GPT-2 did not exist when many ASICs
were designed. Flexibility has enormous value when workloads change.

% -- FLEX: [CORE] The spectrum frames the hardware selection decision
[CORE] Students need this for the cost-performance slide next.
IF AHEAD: ``Where do Apple's Neural Engine and Qualcomm's Hexagon sit?''
IF SHORT: Show the diagram, state the selection rule, move on.
}

% --- Full-width diagram ---
\centering
\safeimg[width=0.88\textwidth,height=4.8cm,keepaspectratio]{accelerator-spectrum.pdf}

\vspace{0.1cm}
\mlsysinsight{Selection rule:}{Match the accelerator's memory system to your workload's access patterns, not its peak FLOPS.}

\end{frame}

\begin{frame}{Heterogeneous SoC and Multi-Chip Scaling}
\note{
% -- LINK: The spectrum showed individual accelerator types; real systems combine them
The flexibility-efficiency spectrum showed CPUs, GPUs, FPGAs, and ASICs as
separate options. Real systems combine multiple accelerators on one chip (SoC)
or across multiple chips (multi-chip modules).

% -- NARRATE: Walk through the SoC block diagram concept
``A modern mobile SoC (Apple M-series, Qualcomm Snapdragon) integrates CPU,
GPU, NPU, and memory on one die. The scheduler routes each layer to the
best accelerator.'' Then multi-chip: ``H100 uses compute die + HBM3 stacks.
TPUv4 pods connect 4,096 chips. Multi-chip scales beyond single-die limits.''

% -- ENGAGE: Why not make one giant chip?
Ask: ``Why not make one enormous chip instead of connecting multiple?''
Give 20 seconds. Expected: yield --- larger dies have more defects.
Cost scales super-linearly with die area.

% -- WARN: Students assume one accelerator handles everything
Common error: ``My model runs on the GPU.'' Correct framing: on mobile
SoCs, different layers run on different accelerators.

% -- FLEX: [CORE] SoC heterogeneity is the reality of edge/mobile deployment
[CORE] Essential for understanding real deployment targets.
IF AHEAD: ``How does chiplet packaging (AMD Zen 4) change the economics?''
IF SHORT: State the SoC concept and the yield argument for multi-chip.
}

\scriptsize
\begin{columns}[T]
  \begin{column}{0.52\textwidth}
    \textbf{System-on-Chip (SoC):} Multiple accelerators, one die.

    \vspace{0.05cm}
    \renewcommand{\arraystretch}{1.0}
    \begin{tabular}{@{}llp{3cm}@{}}
      \toprule
      \textbf{Unit} & \textbf{Strength} & \textbf{Routes To} \\
      \midrule
      CPU  & Serial, branchy logic  & Data preprocessing \\
      GPU  & Parallel graphics/ML   & Conv, MatMul \\
      NPU  & Fixed-function ML      & Quantized inference \\
      DSP  & Signal processing      & Audio, sensor fusion \\
      \bottomrule
    \end{tabular}

    \vspace{0.05cm}
    \begin{mlsyscard}{routingstroke}
    {\scriptsize \textbf{SoC scheduler} routes each layer to the best unit. Performance depends on scheduling quality.}
    \end{mlsyscard}
  \end{column}
  \begin{column}{0.45\textwidth}
    \textbf{Multi-Chip Scaling:}

    \vspace{0.05cm}
    \begin{mlsyscard}{computestroke}
    {\scriptsize
    \textbf{Why multi-chip?}\\
    $\bullet$ Yield: 2$\times$ die area $\approx$ 4$\times$ cost\\
    $\bullet$ H100: compute die + 6 HBM3 stacks\\
    $\bullet$ TPUv4 pod: 4,096 chips\\
    $\bullet$ Chiplets: AMD Zen 4 uses 13 dies\\[0.05cm]
    Multi-chip scales beyond single-die limits.
    }
    \end{mlsyscard}

    \vspace{0.05cm}
    \mlsysconcept{Key:}{SoCs are heterogeneous. Multi-chip is how you scale.}
  \end{column}
\end{columns}

\end{frame}

\begin{frame}{Your Turn: Roofline Diagnosis for a Conv Layer}
\note{
% -- LINK: Roofline was taught abstractly; now apply it to a specific ResNet conv layer
Students learned the Roofline Model. This exercise forces them to calculate
arithmetic intensity for a real convolutional layer and diagnose the regime.

% -- NARRATE: Present problem, then reveal after pause
``ResNet-50 conv3\_1: 3.8 GFLOPs. H100: 989 TFLOPS, 3.35 TB/s. Ridge = 295.
Calculate AI. Compute-bound or memory-bound?'' Wait 90 seconds.

% -- ENGAGE: Peer instruction protocol
Present (30s). Individual work (60s). Key insight: large conv layers with
high channel counts are compute-bound --- unlike the dense layer at batch=1.

% -- WARN: Students may assume all layers are memory-bound
Common error: generalizing from the batch=1 dense example. Correct framing:
conv layers with high channel counts reuse weights extensively, pushing
AI well above the ridge.

% -- FLEX: [CORE] Demonstrates different layers hit different ceilings
[CORE] Shows that one model spans both regimes.
IF AHEAD: ``What happens to AI if you apply 50\% channel pruning (Ch10)?''
IF SHORT: Show solution immediately.
}

\footnotesize
\begin{columns}[T]
  \begin{column}{0.55\textwidth}
    {\normalsize\bfseries Diagnose a Conv Layer}

    \vspace{0.1cm}
    ResNet-50 \texttt{conv3\_1}: \textbf{3.8 GFLOPs}\\
    Input: $28\!\times\!28\!\times\!256$, Output: $14\!\times\!14\!\times\!512$\\
    Weights: $512\!\times\!256\!\times\!3\!\times\!3$ (FP16)

    \vspace{0.1cm}
    H100: 989 TFLOPS, 3.35 TB/s, ridge = 295

    \vspace{0.1cm}
    \textbf{Calculate:}
    \begin{enumerate}\setlength\itemsep{-1pt}
      \item Total bytes transferred
      \item Arithmetic intensity
      \item Compute-bound or memory-bound?
    \end{enumerate}

    {\scriptsize\textcolor{midgray}{(90 seconds --- compare with a neighbor)}}
  \end{column}
  \begin{column}{0.42\textwidth}
    \pause
    \begin{mlsyscard}{datastroke}
    {\scriptsize \textbf{Solution:}\\[0.05cm]
    Weights: $512\!\times\!256\!\times\!9\!\times\!2$ = 2.36 MB\\
    Input: $28^2\!\times\!256\!\times\!2$ = 401 KB\\
    Output: $14^2\!\times\!512\!\times\!2$ = 200 KB\\
    Total $\approx$ \textbf{3.0 MB}\\[0.05cm]
    AI = $3.8\text{G} / 3.0\text{M}$ = \textbf{1,267}\\[0.05cm]
    Ridge = 295. \textcolor{computestroke}{\textbf{Compute-bound!}}\\[0.05cm]
    Contrast: Dense at batch=1 had AI=1\\
    (memory-bound). Same model, different.
    }
    \end{mlsyscard}
  \end{column}
\end{columns}

\end{frame}

\begin{frame}{Cost-Performance: The Right Question}
\note{
% -- LINK: The spectrum showed hardware options; this slide shows how to choose
The flexibility-efficiency spectrum was qualitative. This slide makes the
selection quantitative: performance per dollar, not peak FLOPS.

% -- NARRATE: Walk through the table, then the decision framework
``LLM training: H100 (compute-bound). Batch inference: A100 (throughput).
Real-time batch=1: T4 (memory-bound). On-device: NPU (power budget).''
Key pitfall: ``A100 at batch=1 costs 3--4x more than T4 for identical latency
--- both are memory-bound!'' Decision framework: ``Calculate AI, compare to
ridge point, factor in cost per TFLOP-hour.''

% -- ENGAGE: Test the framework on a real scenario
Ask: ``Your startup serves BERT at batch=1. H100, A100, or T4?'' Expected:
T4 --- all three are memory-bound at batch=1, T4 is cheapest.

% -- WARN: Students default to the most expensive GPU
Common error: ``We need H100 for everything.'' Correct framing: for
memory-bound inference, you are paying for TFLOPS you cannot use.

% -- FLEX: [CORE] Cost-performance reasoning is the practical takeaway
[CORE] Students will use this framework in real hardware selection.
IF AHEAD: ``When does it make sense to use CPUs for ML inference?''
IF SHORT: Cover the table and the T4 vs A100 pitfall.
}

\small
\begin{columns}[T]
  \begin{column}{0.48\textwidth}
    \textbf{Not ``which is fastest?'' but}\\
    \textbf{``which gives best perf/\$?''}

    \vspace{0.15cm}
    \renewcommand{\arraystretch}{1.15}
    {\scriptsize
    \begin{tabular}{@{}llr@{}}
      \toprule
      \textbf{Workload} & \textbf{Best Choice} & \textbf{Why} \\
      \midrule
      LLM training     & H100 & Compute-bound \\
      Batch inference   & A100 & High throughput \\
      Real-time b=1    & T4   & Memory-bound \\
      On-device         & NPU  & Power budget \\
      \bottomrule
    \end{tabular}
    }

    \vspace{0.15cm}
    {\footnotesize\mlsysalert{Pitfall:}{A100 at batch=1 costs 3--4$\times$ more than T4 for same latency.}}
  \end{column}
  \begin{column}{0.48\textwidth}
    \begin{mlsyscard}{crimson}
    {\scriptsize \textbf{Decision framework:}\\[0.05cm]
    1.\ Calculate workload AI\\
    2.\ Compare to ridge point\\
    3.\ AI $<$ ridge $\to$ BW-rich GPU\\
    4.\ AI $>$ ridge $\to$ invest in TFLOPS\\
    5.\ Factor cost per TFLOP-hour
    }
    \end{mlsyscard}
  \end{column}
\end{columns}

\end{frame}

% =============================================================================
\section{Compilers \& Runtime}
% =============================================================================

\begin{frame}{The ML Compilation Pipeline}
\note{
% -- LINK: Fusion and tiling were individual optimizations; the compiler orchestrates all of them
Students saw fusion and tiling as standalone techniques. The ML compiler
combines them all into an optimized execution plan.

% -- NARRATE: Walk through four stages, then the compiler list
``1. Graph optimization: fuse ops, eliminate dead code. 2. Kernel selection:
pick best implementation per op (cuBLAS, cuDNN, auto-tuned). 3. Memory
planning: buffer reuse, in-place ops, liveness analysis. 4. Computation
scheduling: overlap data movement with compute.'' Then the cards: ``Key
compilers: XLA, TorchInductor, TVM, TensorRT. All implement fusion, tiling,
memory planning.''

% -- ENGAGE: Why are there multiple compilers?
Ask: ``Why not one universal ML compiler?'' Expected: different compilers
optimize for different targets (NVIDIA vs AMD vs TPU), different frameworks
(PyTorch vs JAX), and different deployment scenarios (training vs inference).

% -- WARN: Students think they need to implement optimizations manually
Common error: writing custom CUDA kernels for every operation. Correct
framing: compilers handle 80--90\% of optimization automatically. Manual
kernels are only for the critical 10\%.

% -- FLEX: [OPTIONAL] The compiler list is reference material
[OPTIONAL] Four stages are the key concept.
IF AHEAD: ``How does torch.compile work under the hood?''
IF SHORT: Cover the four stages, skip the compiler cards.
}

\scriptsize
\begin{columns}[T]
  \begin{column}{0.55\textwidth}
    \textbf{Four compilation stages:}

    \vspace{0.1cm}
    \textbf{1. Graph Optimization}\\
    Operator fusion, constant folding, dead code elimination.

    \vspace{0.05cm}
    \textbf{2. Kernel Selection}\\
    Pick best implementation per op: cuBLAS, cuDNN, or auto-tuned.

    \vspace{0.05cm}
    \textbf{3. Memory Planning}\\
    Buffer reuse, in-place ops, liveness analysis.

    \vspace{0.05cm}
    \textbf{4. Computation Scheduling}\\
    Overlap data movement with compute; pipeline to hide latency.
  \end{column}
  \begin{column}{0.42\textwidth}
    \begin{mlsyscard}{computestroke}
    {\scriptsize \textbf{Key ML compilers:}\\[0.05cm]
    \textbf{XLA} (TensorFlow/JAX)\\
    \textbf{TorchInductor} (PyTorch 2.0)\\
    \textbf{TVM} (Apache, portable)\\
    \textbf{TensorRT} (NVIDIA inference)\\[0.05cm]
    All implement fusion, tiling, memory planning.
    }
    \end{mlsyscard}

    \vspace{0.05cm}
    \begin{mlsyscard}{routingstroke}
    {\scriptsize \textbf{Runtime:} Dynamic kernel dispatch, memory pools, multi-stream execution.}
    \end{mlsyscard}
  \end{column}
\end{columns}

\end{frame}

% =============================================================================
\section{Wrap-Up}
% =============================================================================

\begin{frame}{Fallacies}
\note{
% -- LINK: Students learned the Roofline; fallacies test if they internalized it
The Roofline and AI framework were taught. These fallacies are common
real-world violations of the principles.

% -- NARRATE: Walk through four fallacies with numbers
First (most common): ``Peak FLOPS: A100 advertises 312 TF. Training: 120--180
TF. Rec models: 10--30 TF. Real-world is 2--3x below peak.'' Second:
``Specialization: Softmax gets 3\% util on A100 but 80--90\% of CPU peak.''
Third: ``Multi-GPU: 8 GPUs give 6--7x, not 8x.'' Fourth: ``A100 and T4
both achieve 4 TF at batch=1 --- same latency, 3--4x price difference.''

% -- WARN: The peak FLOPS fallacy is the most damaging in industry
Common error: comparing GPUs by advertised TFLOPS. Correct framing:
utilization is dictated by arithmetic intensity, not peak throughput.

% -- FLEX: [CORE] Fallacies synthesize the chapter's quantitative reasoning
[CORE] Each fallacy references a specific concept from the lecture.
IF AHEAD: ``How would you benchmark to avoid falling for these?''
IF SHORT: Cover the peak FLOPS and batch=1 fallacies only.
}

\footnotesize
\textbf{Fallacy:} \textit{More specialized hardware always provides better performance.}\\
{\scriptsize Softmax (AI=2--5) achieves 3\% utilization on A100 but 80--90\% of a CPU's lower peak. Specialization helps only when workload matches architecture.}

\vspace{0.05cm}
\textbf{Fallacy:} \textit{Peak FLOPS determine real-world performance.}\\
{\scriptsize A100: 312 TF advertised. Training: 120--180 TF (40--60\%). Rec.\ models: 10--30 TF (3--10\%).}

\vspace{0.05cm}
\textbf{Fallacy:} \textit{Acceleration scales linearly with more GPUs.}\\
{\scriptsize 8-GPU setups achieve 6--7$\times$ (75--87\%) due to sync overhead. Small models: 3--4$\times$ on 8 GPUs.}

\vspace{0.05cm}
\textbf{Fallacy:} \textit{Faster hardware always speeds up inference.}\\
{\scriptsize At batch=1, A100 and T4 both achieve $\sim$4 TF (memory-bound). A100 costs 3--4$\times$ more.}

\end{frame}

\begin{frame}{Pitfalls}
\note{
% -- LINK: Fallacies were conceptual; pitfalls are operational mistakes
Fallacies were wrong mental models. Pitfalls are costly engineering decisions
that teams actually make.

% -- NARRATE: Walk through three pitfalls
First (most damaging): ``300 TFLOPS, 2 TB/s, ridge 150. LayerNorm (AI=1.5)
achieves 3 TFLOPS --- 1\% utilization. A cheaper bandwidth-optimized chip
performs identically.'' Second: ``Dense layer batch=1: A100 and T4 both at
4 TFLOPS. T4 costs 3--4x less.'' Third: ``50+ CUDA kernels = 6--12
engineer-months to port. Isolate behind abstraction layers.''

% -- WARN: The bandwidth pitfall causes the most wasted spend
Common error: buying expensive GPUs for memory-bound workloads.
Correct framing: calculate the workload's AI first, then select hardware.

% -- FLEX: [CORE] Pitfalls connect directly to real procurement decisions
[CORE] These are mistakes students will encounter in their first jobs.
IF AHEAD: ``How do you build hardware abstraction layers in practice?''
IF SHORT: Cover the first pitfall (bandwidth) only.
}

\small
\textbf{Pitfall:} \textit{Ignoring memory bandwidth when selecting accelerators.}\\
{\footnotesize An accelerator with 300 TFLOPS and 2 TB/s has ridge point 150. LayerNorm (AI=1.5) achieves 3 TFLOPS (1\% utilization). A cheaper, bandwidth-optimized alternative performs identically.}

\vspace{0.15cm}
\textbf{Pitfall:} \textit{Deploying small-batch inference on high-compute GPUs.}\\
{\footnotesize Dense layer at batch=1: AI $\approx$ 1 FLOP/byte. Both A100 and T4 achieve $\sim$4 TFLOPS. The T4 costs 3--4$\times$ less.}

\vspace{0.15cm}
\textbf{Pitfall:} \textit{Vendor-specific optimization without portability planning.}\\
{\footnotesize A codebase with 50+ hand-written CUDA kernels requires 6--12 engineer-months to port. Vendor-specific gains (20--40\%) should be isolated behind hardware abstraction layers.}

\end{frame}


\begin{frame}{Muddiest Point}
\note{
% -- NARRATE: Collect muddiest points for next-lecture opening
``Write the single most confusing concept. One sentence.'' Use to open
Ch12 (Benchmarking) with targeted clarification.

% -- FLEX: [CORE] Do not skip even if short
[CORE] 60 seconds, irreplaceable feedback.
IF SHORT: Reduce to 30 seconds, still collect.
}

\centering
\Large\bfseries One-minute paper\\[0.5cm]
\normalsize Write down the \textbf{muddiest point} from today's lecture.\\[0.2cm]
{\small What concept was most confusing or unclear?}\\[0.5cm]
{\footnotesize\textcolor{midgray}{(Hand in on your way out --- or submit digitally)}}
\end{frame}

% --- RETRIEVAL PRACTICE ---
\begin{frame}{What Were the Key Ideas?}
\note{
% -- NARRATE: Retrieval practice before revealing takeaways
``Close notes. Write the 4 most important concepts. 90 seconds.'' Walk around.
Do NOT show Key Takeaways until time is up.

% -- FLEX: [CORE] Highest-impact learning activity
[CORE] Even 60 seconds doubles long-term retention.
IF SHORT: Reduce to 60 seconds but still require writing.
}

\centering
\vspace{1.5cm}
{\Large\bfseries Close your notes.}

\vspace{0.8cm}
{\large Write down the \textbf{4 most important concepts} from today.}

\vspace{0.8cm}
{\normalsize\textcolor{midgray}{90 seconds --- no peeking.}}

\end{frame}

\begin{frame}{Key Takeaways}
\note{
% -- NARRATE: Reveal after retrieval, anchor with numbers
Walk through bullets. ``Memory Wall: 20,000x. Roofline: most ops memory-bound.
Ridge doubles each generation. Systolic arrays: 256x reuse. Fusion: 47 to 8 ms.
Hardware selection: match AI to ridge point. Amdahl: serial caps everything.''

% -- FLEX: [CORE] Definitive summary
[CORE] Students compare retrieval answers against these bullets.
IF SHORT: Read bold terms only.
}

\scriptsize
\begin{itemize}\setlength\itemsep{0pt}
  \item \textbf{Memory Wall}: Moving data costs 20,000$\times$ more energy than computing. Design is about \emph{data locality}.
  \item \textbf{Roofline Model}: Arithmetic intensity determines compute-bound vs.\ memory-bound. Most ML ops are memory-bound.
  \item \textbf{Rising Ridge}: Compute grows $\sim$2$\times$/yr, BW $\sim$1.2$\times$/yr. Ridge points double each generation.
  \item \textbf{Systolic Arrays}: Data reuse via nearest-neighbor PEs. TPUv1: weight loaded once, used 256 times.
  \item \textbf{Kernel Fusion}: Eliminates intermediate HBM writes. ResNet-50: 47$\to$8 ms ($\sim$6$\times$) from fusion alone.
  \item \textbf{Hardware Selection}: Match workload AI to ridge point. Peak FLOPS irrelevant for memory-bound work.
  \item \textbf{Amdahl's Ceiling}: Serial bottlenecks cap acceleration. $S\!=\!\infty$ gives only $1/(1\!-\!p)$ speedup.
\end{itemize}

\end{frame}

\begin{frame}{References}
\note{
% -- NARRATE: Point to canonical papers
``Williams for the Roofline, Horowitz for energy numbers, Jouppi for TPU,
Sze for the survey, H&P for computer architecture foundations.''

% -- FLEX: [OPTIONAL] Reference slide
[OPTIONAL] Not a teaching moment.
IF SHORT: Say ``references on this slide'' and advance.
}

\small
\mlsysref{Williams+09}{Williams, Waterman, Patterson. ``Roofline: An Insightful Visual Model.'' Comm.\ ACM, 2009.}
\mlsysref{Horowitz14}{M. Horowitz. ``Computing's Energy Problem (and what we can do about it).'' ISSCC 2014.}
\mlsysref{Jouppi+17}{Jouppi et al. ``In-Datacenter Performance Analysis of a Tensor Processing Unit.'' ISCA 2017.}
\mlsysref{Sze+17}{Sze et al. ``Efficient Processing of Deep Neural Networks: A Tutorial and Survey.'' IEEE 2017.}
\mlsysref{H\&P21}{Hennessy \& Patterson. \textit{Comp.\ Arch.: A Quantitative Approach}. 6th ed.}
\mlsysref{Gholami+24}{Gholami et al. ``AI and Memory Wall.'' arXiv 2024.}

\end{frame}

\begin{frame}{Next Lecture: Benchmarking}
\note{
% -- LINK: This chapter optimized hardware; next validates the optimizations
We optimized the full D-A-M stack: data selection, model compression,
hardware acceleration. But optimization without measurement is guesswork.

% -- NARRATE: Build anticipation for Ch12
``We claimed speedups from compression and hardware. How do we verify those
claims? Ch12 teaches the Roofline as a diagnostic tool, MLPerf as an
industry standard, and statistical methods to avoid fooling ourselves.''

% -- FLEX: [CORE] Forward hook motivates Ch12 attendance
[CORE] End on: ``Optimization without measurement is guesswork.''
IF SHORT: State the forward hook and dismiss.
}

\small
\begin{columns}[c]
  \begin{column}{0.30\textwidth}
    \centering
    \begin{mlsyscard}{computestroke}
    {\large\bfseries Measure}\\
    {\footnotesize Theoretical FLOPS $\to$\\measured latency}
    \end{mlsyscard}
  \end{column}
  \begin{column}{0.30\textwidth}
    \centering
    \begin{mlsyscard}{datastroke}
    {\large\bfseries Validate}\\
    {\footnotesize Are our optimizations\\real or illusory?}
    \end{mlsyscard}
  \end{column}
  \begin{column}{0.30\textwidth}
    \centering
    \begin{mlsyscard}{routingstroke}
    {\large\bfseries Compare}\\
    {\footnotesize Fair comparison\\across hardware}
    \end{mlsyscard}
  \end{column}
\end{columns}

\vspace{0.3cm}
\centering
{\normalsize Optimization without measurement is guesswork.}\\[0.1cm]
{\small\textbf{From theoretical FLOPs to measured latency --- the Roofline comes to life.}}

\end{frame}


% =============================================================================
% BACKUP SLIDES
% =============================================================================
\appendix

\begin{frame}{Backup: Roofline Model Extended Example}
\note{
% -- NARRATE: Extended roofline calculation for Softmax
Softmax on A100: AI = 3 FLOPs/byte. Ridge = 156. Achievable: 3 x 2,039 GB/s
= 6.1 TFLOPS. Peak: 312 TFLOPS. Utilization: 2\%. Softmax is 78x below
the ridge --- no amount of compute helps.

% -- FLEX: [OPTIONAL] Deploy when students want another calculation
[OPTIONAL] Backup slide.
}
\small
Use if students want another roofline calculation.\\small Softmax on A100: AI $\approx$ 3 FLOPs/byte. Ridge = 156.\\ Achievable performance: $3 \times 2{,}039$ GB/s = 6.1 TFLOPS.\\ Peak: 312 TFLOPS. Utilization: 6.1/312 = \textbf{2\%}.\\ Softmax is 78$\times$ below the ridge --- no amount of compute helps.
\end{frame}

\begin{frame}{Backup: SoC Accelerator Selection Heuristics}
\note{
% -- NARRATE: Practical heuristics for choosing which SoC unit to use
Rule of thumb for mobile SoC scheduling: Conv/MatMul with AI > 50: NPU.
Element-wise ops (ReLU, Add): GPU shader. Control flow and preprocessing:
CPU. Audio/signal: DSP. When in doubt: profile on-device with vendor tools
(Qualcomm Profiler, Apple Instruments, Android GPU Inspector).

% -- FLEX: [OPTIONAL] Deploy during SoC discussion
[OPTIONAL] Backup for heterogeneous deployment questions.
}
\small
Use if students ask how to choose which SoC unit runs each layer.\\[0.1cm]
{\footnotesize
\textbf{Conv/MatMul (AI > 50):} NPU --- fixed-function, highest efficiency.\\
\textbf{Element-wise (ReLU, Add):} GPU shader --- parallel but low AI.\\
\textbf{Control flow, preprocessing:} CPU --- branchy, sequential logic.\\
\textbf{Audio/signal:} DSP --- specialized fixed-point pipelines.\\[0.05cm]
\textbf{Profile tools:} Qualcomm Profiler, Apple Instruments, Android GPU Inspector.\\
\textbf{Key:} On-device profiling beats theoretical analysis. Scheduling overhead can negate unit-level gains.
}
\end{frame}

\begin{frame}{Backup: Systolic Array Data Reuse}
\note{
% -- NARRATE: Alternative explanation of systolic array data reuse
ANALOGY: ``Think of a conveyor belt: weights loaded into registers once.
Inputs flow left to right, one per clock. Each weight multiplied by every
input element. 256x256 array: each weight loaded once, used 256 times.
256x reduction in memory bandwidth demand.''

% -- FLEX: [OPTIONAL] Deploy when students struggle with the concept
[OPTIONAL] Use alternative explanation if main slide was unclear.
}
\small
Use if students struggle with the systolic array concept.\\small Think of a conveyor belt: weights are loaded into registers once.\\ Inputs flow through from left to right, one per clock.\\ Each weight is multiplied by every input element.\\ A 256$\times$256 array: each weight loaded once, used 256 times.\\ This is a 256$\times$ reduction in memory bandwidth demand.
\end{frame}

\end{document}
